\subsection{Morfismos de presentación finita y conjuntos construibles.}
\label{subsection:IV.1.8-fr}

\begin{env}[1.8.1]
\label{IV.1.8.1-fr}
Sea $X$ un preesquema; se sabe que todo abierto cuasicompacto en $X$ es
unión finita de abiertos afines, y recíprocamente. Para que un abierto $U$ de
$X$ sea \emph{retrocompacto} en $X$ (\textbf{0}\textsubscript{III}, 9.1.1), es necesario
y suficiente, por tanto, que para todo \emph{abierto afín} $V$ de $X$, $U\cap V$
sea cuasicompacto. Cuando $X$ es \emph{cuasiseparado} (1.2.1), todo abierto
cuasicompacto en $X$ es retrocompacto en $X$ (1.2.7), por lo que todo subconjunto
localmente construible de $X$ es retrocompacto en $X$
(\textbf{0}\textsubscript{III}, 9.1.13); si \emph{además} $X$ es
\emph{cuasicompacto}, hay identidad entre subconjuntos construibles y subconjuntos
localmente construibles en $X$ (\textbf{0}\textsubscript{III}, 9.1.12), y entre
subconjuntos abiertos construibles y subconjuntos abiertos cuasicompactos
(\textbf{0}\textsubscript{III}, 9.1.5).
\end{env}

\begin{proposition}[1.8.2]
\label{IV.1.8.2-fr}
Sean $X$, $Y$ dos preesquemas, $f:X\to Y$ un morfismo. Para todo subconjunto
construible (\textup{resp.} localmente construible) $Z$ de $Y$, $f^{-1}(Z)$ es
un subconjunto construible (\textup{resp.} localmente construible) de $X$.
\end{proposition}

\begin{proof}
Supongamos $Z$ localmente construible; para todo $x\in X$, habrá una
vecindad abierta afín $V$ de $f(x)$ tal que $Z\cap V$ sea construible
en $V$; si la proposición está demostrada para los subconjuntos construibles,
$f^{-1}(Z)\cap f^{-1}(V)$ será construible en $f^{-1}(V)$, por lo que
$f^{-1}(Z)$ será localmente construible. Basta, por consiguiente, considerar
el caso en que $Z$ es construible, y teniendo en cuenta (\textbf{0}\textsubscript{III},
9.1.3), se reduce al caso en que $Z$ es abierto y \emph{retrocompacto en}
$Y$, es decir, tal que la inyección canónica $Z\to Y$ es un morfismo
cuasicompacto; basta entonces ver que $f^{-1}(Z)$ es
\emph{retrocompacto en $X$}, es decir, tal que la inyección canónica
$f^{-1}(Z)\to X$ es un morfismo cuasicompacto; pero esto resulta de
(1.1.2, (iii)), puesto que $f^{-1}(Z)=Z\times_YX$ (\textbf{I}, 4.4.1).
\end{proof}

\begin{lemma}[1.8.3]
\label{IV.1.8.3-fr}
Sean $X$ un preesquema cuasicompacto y cuasiseparado, $Z$ un subconjunto
construible de $X$. Existen entonces un esquema afín $X'$ y un morfismo
de presentación finita $f:X'\to X$ tales que $f(X')=Z$.
\end{lemma}

\begin{proof}
Como $X$ es cuasicompacto, es unión finita de abiertos afines $X_i$
($i\in I$), y como $X$ es cuasiseparado, resulta de (1.2.7) que cada una
de las inmersiones canónicas $g_i:X_i\to X$ es de presentación finita; se
concluye que si $f_i:X'_i\to X_i$ es un morfismo de presentación finita, lo
mismo ocurre con $g_i\circ f_i:X'_i\to X$ (1.6.2), y si $X'$ es el
preesquema suma de los $X'_i$, el morfismo $h:X'\to X$ que coincide con
$g_i\circ f_i$ en cada $X'_i$ es también de presentación finita
(1.6.5). Como $Z\cap X_i$ es construible en $X_i$
(\textbf{0}\textsubscript{III}, 9.1.8), se ve que se reduce a probar el lema
cuando $X$ es \emph{afín} de anillo $A$. Como $Z$ es entonces unión
finita de conjuntos de la forma $U\cap\complement V$, donde $U$ y $V$ son
abiertos cuasicompactos (\textbf{0}\textsubscript{III}, 9.1.3), se ve, considerando
de nuevo una suma adecuada de preesquemas, que se puede reducir al caso en que
$Z=U\cap\complement V$; como además $U$ es unión finita de abiertos
afines, se puede incluso limitar al caso en que $U$ es \emph{afín}. Como $V$
es cuasicompacto, es unión finita de abiertos de la forma $X_{f_i}$, donde
$f_i\in A$; sea $\mathfrak J$ el ideal de $A$ generado por los $f_i$, y
sea $Z'$ el subesquema cerrado afín de $X$ que define (y que es por
construcción de presentación finita sobre $X$); se tiene por definición
$V=\complement Z'$, por tanto $U\cap\complement V=U\cap Z'$. Consideremos el
esquema afín $X'=Z'\times_XU$
\oldpage[IV-1]{239}
y sea $f:X'\to X$ el morfismo estructural; acabamos de ver que
la inmersión canónica $Z'\to X$ es de presentación finita, y lo mismo ocurre
con la inmersión abierta $U\to X$, que es cuasicompacta puesto que $U$ y
$X$ son afines; se concluye, por tanto, de (1.6.2, (iv)) que $f$ es de
presentación finita, y se tiene $f(X')=Z'\cap U$ (\textbf{I}, 3.4.8), lo que termina la
demostración.
\end{proof}

\begin{theorem}[1.8.4]
\label{IV.1.8.4-fr}
\textup{(Chevalley).} --- Sea $f:X\to Y$ un morfismo de presentación
finita. Para todo subconjunto localmente construible $Z$ de $X$, $f(Z)$ es
localmente construible en $Y$.
\end{theorem}

\begin{proof}
Sean $y\in Y$ y $V$ una vecindad abierta afín de $y$; como $f$ es
cuasicompacto y cuasiseparado, lo mismo ocurre con su restricción
$f^{-1}(V)\to V$, por lo que $f^{-1}(V)$ es un preesquema cuasicompacto y
cuasiseparado (1.2.3); como $Z\cap f^{-1}(V)$ es construible en
$f^{-1}(V)$ (1.8.1), se ve que se puede limitar al caso en que $Y$ es afín
y $Z$ construible; $X$ es entonces cuasicompacto y cuasiseparado, por lo que
(1.8.3), existe un morfismo de presentación finita $g:X'\to X$ tal que
$Z=g(X')$; como $f\circ g$ es de presentación finita, se ve que se
reduce al caso en que $Z=X$; es decir, bastará probar el
\end{proof}

\begin{lemma}[1.8.4.1]
\label{IV.1.8.4.1-fr}
Sean $Y$ un esquema afín, $f:X\to Y$ un morfismo cuasicompacto y
localmente de presentación finita; entonces $f(X)$ es un subconjunto construible
de $Y$.
\end{lemma}

\begin{proof}
Como $X$ es cuasicompacto, es unión finita de abiertos afines; se puede,
por tanto, limitar al caso en que $Y=\operatorname{Spec}(A)$,
$X=\operatorname{Spec}(B)$, siendo $B$ una $A$-álgebra de presentación finita
(1.4.6). Ahora bien, se tiene el
\end{proof}

\begin{lemma}[1.8.4.2]
\label{IV.1.8.4.2-fr}
Sean $A_0$ un anillo, $(A_\lambda)_{\lambda\in L}$ un sistema inductivo de
$A_0$-álgebras, $A=\varinjlim A_\lambda$; si $B$ es una $A$-álgebra de
presentación finita, existen un $\lambda$ y una $A_\lambda$-álgebra de
presentación finita $B_\lambda$ tales que $B$ es isomorfa a
$B_\lambda\otimes_{A_\lambda}A$.
\end{lemma}

\begin{proof}
Por hipótesis, $B$ es isomorfa a un álgebra de la forma
$A[T_1,\ldots,T_n]/\mathfrak J$, donde los $T_i$ son indeterminadas y
$\mathfrak J$ un ideal de tipo finito; sea $(F_j)$ ($1\leq j\leq m$) un
sistema de generadores de $\mathfrak J$. Sea
$\varphi_\lambda:A_\lambda\to A$ el homomorfismo canónico. Como $L$ es
filtrante, existe $\lambda$ tal que cada uno de los coeficientes de cada uno de los
polinomios $F_j$ es la imagen por $\varphi_\lambda$ de un elemento de
$A_\lambda$. Es decir, existe un sistema de polinomios
$(F_{j\lambda})_{1\leq j\leq m}$ de $A_\lambda[T_1,\ldots,T_n]$ tales que
$F_j=\varphi_\lambda(F_{j\lambda})$ para $1\leq j\leq m$. Si
$\mathfrak J_\lambda$ es el ideal de $A_\lambda[T_1,\ldots,T_n]$ generado
por los $F_{j\lambda}$, $\mathfrak J$ es la imagen de
$\mathfrak J_\lambda\otimes_{A_\lambda}A$ en
$A[T_1,\ldots,T_n]=A_\lambda[T_1,\ldots,T_n]\otimes_{A_\lambda}A$;
el anillo $B_\lambda=A_\lambda[T_1,\ldots,T_n]/\mathfrak J_\lambda$ responde a
la cuestión.
\end{proof}

Se aplicará aquí este lema considerando $A$ como límite inductivo de sus
\emph{sub-$\mathbf Z$-álgebras de tipo finito}. Se ve, por tanto, que existe una
tal sub-$\mathbf Z$-álgebra $A_0$ de $A$, y una $A_0$-álgebra de
presentación finita $B_0$ tales que $B$ es isomorfa a
$B_0\otimes_{A_0}A$; pongamos $X_0=\operatorname{Spec}(B_0)$,
$Y_0=\operatorname{Spec}(A_0)$, de modo que se tiene $X=X_0\times_{Y_0}Y$,
siendo $f$ la proyección $X\to Y$; sean $f_0:X_0\to Y_0$,
$g_0:Y\to Y_0$ los morfismos estructurales; resulta de (\textbf{I}, 3.4.8) que
se tiene $f(X)=g_0^{-1}(f_0(X_0))$; teniendo en cuenta (1.8.2), basta, por tanto,
mostrar que $f_0(X_0)$ es construible. Es decir, se reduce finalmente
a demostrar el

\begin{corollary}[1.8.5]
\label{IV.1.8.5-fr}
Sean $Y$ un preesquema noetheriano, $f:X\to Y$ un morfismo de tipo finito.
Para todo subconjunto construible $Z$ de $X$, $f(Z)$ es un subconjunto
construible de $Y$.
\end{corollary}

\oldpage[IV-1]{240}

\begin{proof}
Se reduce, como anteriormente, a probar que $f(X)$ es construible.
Apliquemos el criterio (\textbf{0}\textsubscript{III}, 9.2.3): hay que probar que para
todo subconjunto cerrado irreducible $T$ de $Y$, $T\cap f(X)=f(f^{-1}(T))$ es
raro en $T$ o contiene un abierto no vacío de $T$; tomando un
subpreesquema de $Y$ que tenga a $T$ por espacio subyacente (\textbf{I}, 5.2.1) y
considerando su imagen inversa por $f$, se ve finalmente (teniendo en cuenta
(1.5.4, (iii))) que se reduce a probar que, si se supone $Y$
\emph{irreducible} y $f$ \emph{dominante}, $f(X)$ contiene un abierto no vacío de
$Y$. Se puede además suponer $Y$ afín; entonces $X$ es unión finita
de abiertos afines $V_i$, y como $f(X)$ es denso en $Y$, lo mismo ocurre
con al menos uno de los $f(V_i)$. Se puede, por tanto, también suponer $X$ afín; si
$(X_j)$ es la familia (finita) de las componentes irreducibles de $X$, se ve,
como anteriormente, que al menos uno de los $f(X_j)$ es denso en $Y$; se puede, por tanto,
suponer $X$ irreducible. Por último, sustituyendo $f$ por $f_{\mathrm{red}}$
(1.5.4, (vi)) se pueden suponer $X$ y $Y$ \emph{íntegros}. Entonces
$X=\operatorname{Spec}(B)$, $Y=\operatorname{Spec}(A)$, donde $A$ es un
anillo íntegro noetheriano, $B$ una $A$-álgebra íntegra de tipo finito que contiene
a $A$ (\textbf{I}, 1.2.7). Basta mostrar que existe $g\in A$ tal que, para todo
$y\in D(g)$ (es decir, tal que $g\notin\mathfrak j_y$) el ideal
$\mathfrak j_y$ es la intersección de $A$ y de un ideal primo de $B$, pues
esto mostrará que $D(g)\subset f(X)$. Finalmente basta probar el
\end{proof}

\begin{lemma}[1.8.5.1]
\label{IV.1.8.5.1-fr}
Sean $A$ un anillo íntegro, $B$ una $A$-álgebra íntegra de tipo finito.
Existe $g\in A$ tal que todo homomorfismo de $A$ a un cuerpo
algebraicamente cerrado $\Omega$, no nulo en $g$, se prolonga a un homomorfismo
de $B$ a $\Omega$.
\end{lemma}

\begin{proof}
Ahora bien, este es un resultado clásico de álgebra conmutativa (Bourbaki,
\emph{Alg. comm.}, cap. V, \S~3, n\textsuperscript{o}~1, cor. 3 del teorema 1).
\end{proof}

\begin{corollary}[1.8.6]
\label{IV.1.8.6-fr}
Sean $Y$ un preesquema irreducible, $\eta$ su punto genérico,
$f:X\to Y$ un morfismo localmente de tipo finito. Si $f^{-1}(\eta)$ no es
vacío, existe una vecindad abierta $U$ de $\eta$ en $Y$ tal que
$f^{-1}(y)$ es no vacío para todo $y\in U$. Si $\mathscr F$ es un
$\mathcal O_X$-Módulo cuasicoherente de tipo finito, y si
$\mathscr F_\eta=\mathscr F\otimes_{\mathcal O_Y}\boldsymbol{k}(\eta)$ no es nulo,
entonces existe una vecindad abierta $U'$ de $\eta$ en $Y$ tal que
$\mathscr F_y=\mathscr F\otimes_{\mathcal O_Y}\boldsymbol{k}(y)$ no es nulo para
todo $y\in U'$.
\end{corollary}

\begin{proof}
Si $p_y$ es la proyección canónica de la fibra
$f^{-1}(y)=X\times_Y\operatorname{Spec}(\boldsymbol{k}(y))$ en $X$, se tiene
$\mathscr F_y=p_y^*(\mathscr F)$, por tanto
$\operatorname{Supp}(\mathscr F_y)=p_y^{-1}(\operatorname{Supp}(\mathscr F))
=\operatorname{Supp}(\mathscr F)\cap f^{-1}(y)$ en virtud de (\textbf{I}, 9.1.13.1) y
(\textbf{I}, 3.6.1), puesto que $\mathscr F$ es de tipo finito; además
$\operatorname{Supp}(\mathscr F)$ es cerrado en $X$ (\textbf{0}\textsubscript{I},
5.2.2), y si $Z$ es un subpreesquema cerrado de $X$ que tiene por espacio
subyacente $\operatorname{Supp}(\mathscr F)$ (\textbf{I}, 5.2.1), y $j$ la inmersión
canónica $Z\to X$, $f\circ j$ es localmente de tipo finito (1.3.4); esto
prueba que la primera afirmación implica la segunda. Para probar la
primera afirmación, observemos primero que se pueden suponer $X$ y $Y$
reducidos (1.3.4, (vi)), es decir, $Y$ \emph{íntegro}. Sea $\xi$ un punto de
$f^{-1}(\eta)$; se pueden sustituir $X$ y $Y$ por vecindades abiertas
afines de $\xi$ y $\eta$ respectivamente, y por consiguiente suponer que
$X=\operatorname{Spec}(B)$, $Y=\operatorname{Spec}(A)$, siendo $B$ una
$A$-álgebra de tipo finito. Además, si $Z$ es el subpreesquema cerrado reducido
de $X$ que tiene a $\overline{\{\xi\}}$ por conjunto subyacente (\textbf{I}, 5.2.1), se
puede, como anteriormente, sustituir $X$ por $Z$, es decir, suponer $B$
\emph{íntegro} (\textbf{I}, 5.1.4). Como el morfismo $f$ es entonces dominante,
el homomorfismo correspondiente $\varphi:A\to B$ es inyectivo (\textbf{I}, 1.2.7);
por tanto el corolario resulta finalmente del lema (1.8.5.1).
\end{proof}

\oldpage[IV-1]{241}

\begin{proposition}[1.8.7]
\label{IV.1.8.7-fr}
Sean $S$ un preesquema, $g:X\to S$, $h:Y\to S$ dos morfismos de
presentación finita, $f:X\to Y$ un $S$-morfismo. Para todo $s\in S$, pongamos
$X_s=g^{-1}(s)=X\times_S\operatorname{Spec}(\boldsymbol{k}(s))$,
$Y_s=h^{-1}(s)=Y\times_S\operatorname{Spec}(\boldsymbol{k}(s))$,
$f_s=f\times 1:X_s\to Y_s$. Entonces el conjunto de los $s\in S$ tales que $f_s$
es sobreyectivo (\textup{resp.} radicial) es localmente construible.
\end{proposition}

\begin{proof}
Sea $E$ el conjunto de los $s\in S$ tales que $f_s$ es sobreyectivo; el conjunto
$S\setmin E$ es igual a $h(Y\setmin f(X))$; ahora bien, $f$ es de presentación finita (1.6.2, (v))
por lo que $f(X)$ es localmente construible en $Y$ (1.8.4); como $h$ es de
presentación finita, $h(Y\setmin f(X))$ es localmente construible en $S$, y por tanto
también $E$.

Para mostrar que el conjunto $E'$ de los $s\in S$ tales que $f_s$ es radicial
es localmente construible, utilizaremos el
\end{proof}

\begin{lemma}[1.8.7.1]
\EGAReaderTarget{ega4.c04.lem.1-8-7-1}{IV.1.8.7.1-fr}
Sean $U$, $V$ dos preesquemas; para que un morfismo $f:U\to V$ sea
radicial, es necesario y suficiente que el morfismo diagonal
$\Delta_f:U\to U\times_VU$ sea sobreyectivo. Por consiguiente, todo morfismo radicial
es separado.
\end{lemma}

\begin{proof}
En efecto, como $\Delta_f$ es una inmersión (\textbf{I}, 5.3.9), es un morfismo
localmente de tipo finito (1.3.4); decir que $\Delta_f$ es sobreyectivo significa,
por tanto, que para todo cuerpo algebraicamente cerrado $K$, la aplicación
correspondiente $U(K)\to(U\times_VU)(K)=U(K)\times_{V(K)}U(K)$ es sobreyectiva
(1.3.7 y \textbf{I}, 3.4.2.1); pero por definición de un producto fibrado de conjuntos,
esto significa que la aplicación $U(K)\to V(K)$ correspondiente a $f$ es
inyectiva, y esto equivale a decir que $f$ es radicial (\textbf{I}, 3.5.5).

Siendo así, $\Delta_{f_s}$ se deduce de $\Delta_f:X\to X\times_YX$ por el
cambio de base $\operatorname{Spec}(\boldsymbol{k}(s))\to S$ (\textbf{I}, 5.3.4). Como $f$
es de presentación finita, lo mismo ocurre con el morfismo estructural
$X\times_YX\to Y$ (1.6.2, (iv)), y por tanto también con $\Delta_f$ (1.6.2, (v));
basta, pues, aplicar la primera parte de la proposición a $\Delta_f$,
utilizando el lema (1.8.7.1).
\end{proof}

